\documentclass[a4paper,12pt]{article}
\usepackage[spanish,activeacute]{babel}
\usepackage[latin1]{inputenc}

%% Para las columnas
\usepackage{multicol}
\usepackage{amssymb,amsmath,eepic,fancyhdr}
\usepackage{latexsym}
\usepackage{datenumber}


%% Para numeraci?n autom?tica
\newcounter{nroproblema}
\setcounter{nroproblema}{1}
\newcounter{nroapartado}
\setcounter{nroapartado}{1}

\newcommand{\n}{%
    \vspace{.1in}\noindent{\bf \arabic{nroproblema}.\ }%
    \addtocounter{nroproblema}{1}%
    \setcounter{nroapartado}{1}%
    }%

\newcommand{\nn}{%
    \vspace{.1in}\noindent \hskip.4cm {\bf \arabic{nroproblema}.\ }%
    \addtocounter{nroproblema}{1}%
    \setcounter{nroapartado}{1}%
    }%


\newcommand{\ap}{\vspace*{0.2cm}\hspace*{0.2cm}%
   {\bf \alph{nroapartado})}
   \addtocounter{nroapartado}{1}%
   }

 
%marcos by Marco
\newcommand{\pd}[1]{\frac{\partial}{\partial #1}} %\pd{x}
\newcommand{\NN}{\ensuremath{\mathbb N}}
\newcommand{\ZZ}{\ensuremath{\mathbb Z}}
\newcommand{\CC}{\ensuremath{\mathbb C}}
\newcommand{\RR}{\ensuremath{\mathbb R}}
\newcommand{\TT}{\ensuremath{\mathbb T}}
\newcommand{\g}{\ensuremath{\frak{g}}}
\newcommand{\h}{\ensuremath{\frak{h}}}
\newcommand{\ft}{\ensuremath{\frak{t}}}
\newcommand{\vX}{\frak{X}}
\newcommand{\cB}{\mathcal{B}}
\newcommand{\cN}{\mathcal{N}}
\newcommand{\cL}{\mathcal{L}}
\newcommand{\cD}{\mathcal{D}} 
 

%% m�rgenes

% Anterior:

%\setlength{\unitlength}{1mm} \addtolength{\voffset}{-15mm}
%\addtolength{\textheight}{30mm} \addtolength{\hoffset}{-20mm}
%\addtolength{\textwidth}{45mm}

\setlength{\unitlength}{1mm} \addtolength{\voffset}{-22mm}
\addtolength{\textheight}{40mm} \addtolength{\hoffset}{-22mm}
\addtolength{\textwidth}{46mm}




\begin{document}

\noindent
{\sffamily\large Marco Zambon 
\hfill  Master UAM, 2013-14  \\Geometr\'ia simpl\'ectica y de Poisson}


 

\begin{center}\bfseries\large\textsf{Hoja 7    \vspace{-2mm}}
\end{center}
  
 
 

\n  Considera $\RR^3$ y una funci\'on $f\colon \RR\to \RR$. Considera
el campo bivectorial 
$$\pi:=f(z)\pd{x}\wedge\pd{y}.$$

\ap Comprueba que $\pi$ es un campo bivectorial  de Poisson.
 
\ap Cuales son las hojas simpl\'ecticas de $\pi$? Dib\'ujalas. 

  
\n[3 puntos]
Sea $V$ un espacio vectorial y $\pi\in \wedge^2 V$. Denota por $\sharp\colon V^*\to V$ la aplicaci\'on lineal $\xi \mapsto \iota_{\xi}\pi$.
Demuestra que
  $$Ker(\sharp)=[Im(\sharp)]^{\circ}.$$
  
  
%\{,\cdot,\cdot\}_M  
\n Sean $(M,\pi_M)$ y   $(N,\pi_N)$
variedades de Poisson, y $\phi\colon M\to N$ una aplicaci\'on.

Vamos a demostrar en 3 pasos que:

 \emph{$\phi$ es un morfismo de Poisson $\Rightarrow$
$\text{grafo}(\phi)$ es una subvariedad coisotropica de $(M\times N, \pi)$},

donde $\pi$ es el campo bivectorial dado por $\pi|_{(m,n)}=(\pi_M)|_m-
(\pi_N)|_n$ para todo $(m,n)\in M\times N$.
  

\noindent\textbf{Notaci\'on:} Para todo $f\in C^{\infty}(N)$, denotamos $$F:=\phi^*(f)-f  \in 
C^{\infty}(M\times N).$$  (Aqu\'i estoy abusando un poco la notaci\'on, denotando por $f$ una funci\'on en $N$ y tambi\'en la funci\'on en $M\times N$ dada por $(m,n)\mapsto f(n)$. Lo mismo vale para $\phi^*(f)$.)
  
\ap Denota por $I_{  \text{grafo}(\phi)}$ el conjunto de funciones en $M\times N$ que se anulan en $\text{grafo}(\phi)$. Demuestra: para todo $f\in C^{\infty}(N)$, tenemos $$F\in I_{  \text{grafo}(\phi)}.$$ 


\ap Para todo $f_1,f_2\in C^{\infty}(N)$, tenemos $$\{F_1,F_2\}_{M\times N}\in I_{  \text{grafo}(\phi)}.$$

\ap Demuestra que $$\{I_{  \text{grafo}(\phi)},I_{  \text{grafo}(\phi)}\}_{M\times N}\subset I_{  \text{grafo}(\phi)}.$$

\noindent\textbf{Nota:} Puedes utilizar que $\{F:f\in C^{\infty}(N)\}$
genera $I_{  \text{grafo}(\phi)}$ como ideal multiplicativo, es decir, cada elemento de $I_{  \text{grafo}(\phi)}$ es una suma finita  $\sum_i G_iF_i$ donde $G_i\in C^{\infty}(M\times N)$ y   $F_i=\phi^*(f_i)-f_i$ para  $f_i\in C^{\infty}(N)$.

   \end{document}
  \noindent\textbf{Consejo:}
\noindent\textbf{Nota:}